\chapter{总结与展望}\label{chap:conclusion}
\section{总结}
本毕业设计围绕电缆通道检测机器人的运动控制与路径规划问题展开研究，针对复杂狭窄环境下机器人自主巡检能力不足的问题，完成了轮腿式底盘的建模、控制与路径规划方法设计。首先，建立了串联五连杆轮腿式机器人的运动学与动力学模型，并将其等效简化为一阶倒立摆模型，设计了基于 LQR 的状态反馈控制器，再通过 VMC 方法将控制量映射为轮电机与关节电机的执行输入。MATLAB/Simulink 仿真结果表明，该控制方案能够使机器人在平坦地面及包含四个 0.25 m 障碍物的场景下保持姿态稳定，并顺利完成越障运动。

在路径规划方面，本文构建了面向动态障碍环境的强化学习仿真平台，对 DWA、DQN、SAC 以及引入 Transformer 的 TransSAC 等算法进行了分析与调试，并结合任务特点设计了奖励函数与分布式训练框架。仿真结果表明，传统 DWA 算法在复杂场景下容易陷入局部最优，到达率仅为 0.53，综合性能有限；DQN 在经过充分训练后到达率可稳定在 0.9 左右，具有较好的收敛性；TransSAC 则进一步利用时序上下文信息提升了策略学习效率，在复杂动态环境中表现出更快的收敛速度和更优的路径效率。总体而言，本文实现了面向电缆通道检测任务的运动控制与智能导航方法验证，为后续工程化应用提供了参考。
\section{展望}
本文的研究为电缆通道检测机器人的自主运动提供了较为完整的控制与规划框架，但仍存在进一步提升的空间。首先，当前控制与规划主要基于仿真环境验证，后续可结合真实样机与现场测试，对模型参数、控制律以及奖励函数进行针对性修正，以增强方法的工程适配性和鲁棒性。其次，本文的路径规划主要面向二维导航场景，未来可以进一步融合三维环境感知、语义地图与多传感器信息，使机器人对复杂电缆通道中的坡道、台阶、积水和狭窄分支具备更强的识别与决策能力。

此外，TransSAC 虽然在复杂动态环境中表现出较好的路径效率，但其训练过程对超参数和时序窗口较为敏感，后续可尝试引入更轻量的时序建模结构、自适应奖励调整机制以及更高效的分布式训练策略，以降低训练成本并提升泛化能力。随着感知、控制与智能决策技术的持续发展，电缆通道检测机器人有望在自动巡检、缺陷识别和全自主运维等方向取得更广泛的应用。